\section{Sharpening the finite-test prices}
\label{sec:background-test-prices}
The real parameter tests of Section~\ref{sec:background-ray-estimates}
need small prices on one fixed parameter interval. The local expansion
first provides a bound on a growing interval and a better pointwise
bound. A finite polynomial interpolation combines these two estimates.
We give the interpolation and its numerical purchase, keeping the
underlying wide-interval analytic estimate visible.

\subsection{A finite interpolation estimate}
Let $\mathcal E$ be a vector space of families on a fixed support
dictionary, with a nonnegative budget $\mathcal N$ satisfying the
triangle inequality and absolute homogeneity. A budget defined as
an infimum over admissible summable majorants has these properties:
add or scale the majorants and then take their infima. The argument
below equally applies by retaining the majorants themselves.
A finite test is a finitely supported real signed measure on
$[-1,1]$,
\[
 \lambda=\sum_{h=1}^{h_0}a_h\delta_{t_h},\qquad a_h\in\mathbb R,
 \quad t_h\in[-1,1],
\]
and for a function $G$ on $[-1,1]$ with values in a real vector
space we write $\int G\,d\lambda=\sum_ha_hG(t_h)$. Set
\[
 M_r(\lambda)=\sup_{n\ge0}\left|\int(t/r)^n\,d\lambda(t)\right|,
                    \qquad r\ge1.
\]
Let $F:[-1,1]\to\mathcal E$. Suppose
\[
 \mathcal N\left(\int F\,d\lambda\right)\le B M_v(\lambda),
 \qquad \mathcal N(F(t))\le b\quad(-1\le t\le1),
\]
where $1\le q\le v$ and $b,B\ge0$, the first bound being assumed
for every finite test $\lambda$. The integrals here are finite
linear combinations; no measurability of $F$ is required.
For each integer $k\ge0$, a constant $U=U(k,q)\ge0$, independent
of the family, its dictionary and its volume, gives
\begin{equation}\label{eq:background-sharp-test}
 \mathcal N\left(\int F\,d\lambda\right)
 \le M_q(\lambda)\left\{Ub+(1+U)(q/v)^{k+1}B\right\}.
\end{equation}

Choose the $k+1$ distinct points $t_i=i/(k+1)$, $0\le i\le k$,
and their Lagrange polynomials $\ell_i$ of degree at most $k$.
Write $\ell_i(t)=\sum_{j=0}^kc_{ij}t^j$ and put
\[
                    U_i=\sum_{j=0}^k|c_{ij}|q^j,
                    \qquad U=\sum_{i=0}^kU_i.
\]
For $\eta=M_q(\lambda)$, we have
$|\int\ell_i\,d\lambda|\le\eta U_i$. Define
\[
              \Pi\lambda=\sum_{i=0}^k
                         \left(\int\ell_i\,d\lambda\right)\delta_{t_i}.
\]
Lagrange interpolation of each polynomial of degree at most $k$
gives $\int t^j\,d(\lambda-\Pi\lambda)=0$ for $j\le k$.
For $j>k$, use $|t_i/q|\le1$ to get
\[
 \left|\int(t/q)^j\,d(\lambda-\Pi\lambda)\right|
          \le\eta+\sum_i\left|\int\ell_i\,d\lambda\right|
          \le\eta(1+U).
\]
Since $q/v\le1$, multiplication by $(q/v)^j$ proves
\[
 M_v(\lambda-\Pi\lambda)\le\eta(1+U)(q/v)^{k+1}.
\]
The pointwise price and the triangle inequality give
$\mathcal N(\int F\,d\Pi\lambda)\le\eta Ub$.
Adding this to the wide-interval bound on $\lambda-\Pi\lambda$
proves \eqref{eq:background-sharp-test}. Thus the same constant
works for the full regular derivative budget, including all its
site sums; no count of dictionary entries enters the estimate.

\subsection{The fixed-radius price at large coupling}
Fix the stage and block length, and put $y=x+a$, where $a>0$
is the fixed reference coupling increment. The broad analytic
estimate to be sharpened has radius and price
\[
                 v(x)=\tfrac12 y^{1/256},\qquad B(x)=y^{-1/4}.
\]
The local pointwise price has the form
\[
             b(x)=3t C_{\rm st}x^{-11/10}+R(x),
\]
where $t>0$ is a small response parameter and $C_{\rm st}>0$
is the fixed stage conversion constant. The remainder $R$ has
upper exponent at most $-23/20$: concretely, for every
$\epsilon>0$, $|R(x)|\le C_\epsilon y^{-23/20+\epsilon}$
eventually. This estimate comprises the local comparison and
truncation errors; it is separate from the broad analytic bound.
For the selected one-step source it is the regular budget proved in
Section~\ref{sec:selected-nonlinear-offering}. There the principal
and complementary regular prices are purchased as $E/16$ and
$2E/16$, with $E=C_{\rm st}x^{-11/10}$; here they are purchased with
the fraction $t\le1/16$ in place of $1/16$; the smaller regular
reserve is the one permitted in
Section~\ref{sec:common-funded-history} and is chosen before the
entry delays. The offering section states its budget only for the
fraction $1/16$. Its proof uses the principal and complementary
prices only through their sum in the truncation bound
\eqref{eq:selected-offer-truncation}, so with the hypotheses $tE$
and $2tE$ the same proof gives the leading term
$3tC_{\rm st}x^{-11/10}$ in place of $3E/16$. The other terms of
that budget are unchanged: in the notation of that section,
\[
          R=C_{\rm turn}g+g+C_{\rm turn}c'e^{A_cM},
          \qquad g=y^{-23/20},
\]
where $c'$ is the price of the folded counterterms. Both $g$ and
$c'$ have upper exponent $-23/20$, and $C_{\rm turn}$, $A_c$, $M$
are fixed at the stage. The bound holds at every point of the
parameter line, for the raw family in its one common dictionary.

Fix $q\ge2$ and $s>0$, take $k=300$ in
\eqref{eq:background-sharp-test}, and choose $t$ small enough that
$6Ut\le s$. The leading pointwise contribution is at most
$(s/2)C_{\rm st}x^{-11/10}$. The other part of the bound is
\[
 UR(x)+(1+U)(2q)^{301}y^{-301/256-1/4}.
\]
Its upper exponent is at most $-23/20$, because
$301/256+1/4=365/256>23/20$. Since $23/20>11/10$
and $y/x\to1$, that expression is eventually bounded above by
$(s/2)C_{\rm st}x^{-11/10}$. Hence the fixed-radius test price is
\begin{equation}\label{eq:background-pre-reset-price}
                         sC_{\rm st}x^{-11/10}.
\end{equation}
There is no requirement that $U$ be small. It depends only on
$300$ and $q$; the response parameter is chosen after it, and the
coupling threshold after all those choices. For example
$t\le s/[6(1+U)]$ always suffices. Since $v(x)\to\infty$,
the condition $q\le v(x)$ also holds beyond that threshold.

\subsection{The reset and the final regular price}
The reset and local Taylor subtraction are linear in the raw
logarithmic family and in its extracted coupling coefficient.
On a fixed stage, the regular output price is at most
\begin{equation}\label{eq:background-regular-price}
                   C_0 H+C_1|\alpha|\theta(x)^3,
\end{equation}
where $C_0,C_1\ge0$ are independent of the coupling, period and
input family, $H$ is the raw test price, and
$\theta(x)=x^{-1/2}\exp((\log\log x)^2)$ is the fluctuation
threshold. The two constants collect respectively the reset,
low-path and high-path derivative prices, and the high-path
remainder. This is the same local subtraction estimate used in
the nonlinear transition, applied to a finite linear combination
before taking its majorant. It is the regular cost $E_{\rm new}$ of
Section~\ref{sec:physical-reset-reorganization}, with the threshold
of that section equal to $\theta(x)$: in the notation of
\eqref{eq:reset-stage-constants},
\[
 C_0=B^{\rm reset}=C_{\rm reset}+C_{\rm turn}(L_0+L_1),\qquad
 C_1=v^{\rm reset}=C_{\rm turn}L_1.
\]
The estimate is used under the hypotheses of that section, for the
path, clipping data and outgoing budget of the transition under
consideration. The path satisfies its paid bounds at the path width
$\sigma$. The relaxation radius $r_c$ of the clipping data and the
support exponent $b$ and radius $B$ of the budget satisfy
$12\le r_c\sigma$ and $(b+1)/B\le\sigma/4$. The direct clipping
coefficient is at most one, $B\ge2$, and
$2\le\theta(x)e^{(B-2)/r_c}$, as in \eqref{eq:reset-clip-conditions}.
The group is simple. The path correction $\alpha$ is one fixed number along the
whole parameter line, with $|\alpha|\le A$. With these data the
constants $C_0,C_1$ depend only on the stage.

The stage conversion constant is selected as
\[
                 C_{\rm st}=\frac1{64(1+C_0)}.
\]
This is the allocation $c^{\rm reset}$ of
\eqref{eq:reset-stage-constants}.
It is positive and satisfies both $C_{\rm st}\le1/64$ and
$C_0C_{\rm st}\le1/64$. Fix $d>0$, take $0<s\le d$ in
\eqref{eq:background-pre-reset-price}, and retain the fixed
bound $|\alpha|\le A$. The first term of
\eqref{eq:background-regular-price} is at most
$(d/64)x^{-11/10}$. The second is $o(x^{-11/10})$, since
\[
 \frac{\theta(x)^3}{x^{-11/10}}
          =x^{-2/5}\exp(3(\log\log x)^2)\longrightarrow0.
\]
Increasing the coupling threshold therefore gives regular
output price at most $d x^{-11/10}$.

This remains a finite-test statement. If the moment bound of
the finite test is $\eta$, the raw price is $\eta H$, the
constant path correction becomes the same finite combination
of $\alpha$, and its modulus is at most $\eta A$ by the zeroth
moment test. The linear reset and the triangle inequality thus
multiply \eqref{eq:background-regular-price} by $\eta$.
Dividing by the positive regular normalization gives the regular
part of \eqref{eq:background-ray-tests}.
The scalar coupling coefficient is a bounded linear projection
of the raw family, with price $H/\theta(x)^2$. Consequently its
finite-test price is at most the scalar allowance
$f_m=sC_{\rm st}x^{-11/10}/\theta(x)^2$ used there.

\subsection{The forced price and the remaining analytic input}
A finite-test bound alone may choose its majorant after the test,
so recovery of a power series for each value separately does not yet
produce one common support majorant. We supply that step. Let $S_i$
be the countable list of nonempty supports and let $e\ge0$ be the
cover exponent. Suppose that for every finite test $\lambda$ there
are numbers $k_\lambda(i)\ge0$ such that
\[
 \begin{gathered}
 \left|\int F(t,i,U)\,d\lambda(t)\right|\le k_\lambda(i)
 \quad\hbox{for every configuration }U,\\
 \sum_{i:z\in S_i}k_\lambda(i)e^{e C(S_i)}
                         \le B M_q(\lambda)
 \quad\hbox{at every root }z.
 \end{gathered}
\]
Here $C(S_i)$ is the cover number. Nonemptiness and $e\ge0$
imply $k_\lambda(i)\le B M_q(\lambda)$ for each $i$.
For each entry and configuration, the closed-span argument of
Section~\ref{sec:background-ray-estimates} therefore extends
finite-test evaluation to a bounded functional $T_{i,U}$ on the
closed span of the moment vectors $v(t)=((t/q)^n)_n$.

Fix $n$. Choose finite tests $\lambda_j$ whose moment vectors
$v(\lambda_j)$ converge in sup norm to $e_n$; these approximants
exist by the coordinate-vector construction above. In particular
$M_q(\lambda_j)\to1$. Choose their majorants $k_j$ as in the
hypothesis. Their coordinates lie in one common bounded interval.
A diagonal subsequence on the countable indices $i$ gives
$k_j(i)\to b_{n,i}\ge0$. For each finite set of indices and
root $z$, pass to the limit in its weighted sum. Taking the
supremum over finite sets yields
\[
                 \sum_{i:z\in S_i}b_{n,i}e^{e C(S_i)}\le B.
\]
Also $|T_{i,U}e_n|\le b_{n,i}$, by continuity of $T_{i,U}$.
This last conclusion holds for every $U$ on the same subsequence:
the chosen majorants bounded all configurations from the start,
and convergence of the moment vectors is independent of $U$.
No subsequence indexed by configurations is selected.

For $1\le r<q$, put
\[
             K_i=\sum_{n\ge0}(r/q)^n b_{n,i}.
\]
The bounds $b_{n,i}\le B$ make each series finite. Tonelli's
theorem for the nonnegative weighted series gives
\[
                 \sum_{i:z\in S_i}K_i e^{e C(S_i)}
                              \le\frac{B}{1-r/q}.
\]
The coefficients $(r/q)^nT_{i,U}e_n$ have total absolute sum
at most $K_i$ and reconstruct $F(t,i,U)$ as a series in $t/r$.
Thus one common support majorant works for every configuration.
Applied to the soft correction from the linear reset with
$r=q/2$, this proves the required factor-two price and radius loss.

The compatible
finite-subset products preserve absolute coefficient summability:
their coefficients obey convolution, and the sum of coefficient
norms of a product is bounded by the product of the corresponding
sums. The proof is Tonelli's theorem for the nonnegative coefficient
majorants. Thus the forced absorption argument of
Section~\ref{sec:soft-log-fusion} applies also to
these parameter coefficient prices, with the same enlarged
support reserve and the same scalar parameter $\kappa$.
Its inputs are the following. The hard price $a_{\rm h}$ is the
common coefficient majorant of the raw forced family at radius $q/2$
and support exponent $b+1$, where $b$ is the outgoing support
exponent. The soft price is any $c_{\rm s}\ge2S$, where $S$ is the
soft cost $S_{\rm new}$ of
Section~\ref{sec:physical-reset-reorganization} at the raw price $H$
and the fixed correction $\alpha$; the factor two is the radius loss
just proved. If $K_m>0$ denotes the prescribed forced normalization
and $\kappa\ge0$ satisfies
\begin{equation}\label{eq:background-forced-payment}
 c_{\rm s}\le1,\qquad (2B+1)^4I\le1,\qquad
 e^{-\kappa}(e^I-1)\le dK_m,\qquad I=e^\kappa a_{\rm h}+2c_{\rm s},
\end{equation}
then the final forced array has one common coefficient majorant at
radius $q/2$ and exponent $b$ with rooted price at most $dK_m$.
The existence of such a $\kappa$ at every sufficiently late entry,
uniformly over the coupling band, is a numerical input of this
section and is not derived here. It is a payment of the same kind
as \eqref{eq:nonlinear-payment}, but it is not a consequence of that
statement alone: the hard price $a_{\rm h}$ is the forecast
allowance reached in Section~\ref{sec:projected-parameter-prices},
and the payment uses the purchased gap between the exponential rate
of that forecast and the smaller rate of the output normalization
$K_m$.
The test inequality follows by applying each finite moment test
to the absolutely summable coefficient series. This supplies the
forced clause of \eqref{eq:background-ray-tests} whenever the
broad ordinary and forced analytic estimates have been supplied.

The broad estimates used in this section are substantive: the
ordinary raw logarithmic family must have radius $y^{1/256}/2$
and price $y^{-1/4}$ in the full derivative budget, and the raw
forced family must have its common coefficient majorant at the
paid broad radius. The pointwise remainder $R$ must have the
stated upper exponent. The interpolation and reset estimates
above do not themselves prove those statements. They explain
exactly how those local analytic estimates, once obtained for
the same transition and parameter line, give the small prices
needed by the coupled history argument.
